\chapter{Código fuente del simulador Arduino}
\label{anexo:arduino}

Este anexo recoge los fragmentos más representativos del código fuente del simulador
de ECU implementado sobre Arduino MKR. El código completo está disponible en el
repositorio del proyecto (\texttt{arduino-sim/ecu\_sim/}).

\section*{A.1 — Definiciones de protocolo (\texttt{ecu\_protocol.h})}

El fichero de cabecera centraliza las constantes del protocolo, las macros de
configuración y las funciones de codificación OBD-II.

\begin{lstlisting}[language=C++, caption={Constantes de CAN, ISO-TP, UDS y configuración — ecu\_protocol.h}]
// ---------- CAN BUS ----------
#define CAN_SPEED          500E3
#define ECU_CAN_ID         0x7E0
#define ECU_RESPONSE_ID    0x7E8

// ---------- ISO-TP ----------
#define ISOTP_PCI_SF       0x00  // Single Frame
#define ISOTP_PCI_FF       0x10  // First Frame
#define ISOTP_PCI_CF       0x20  // Consecutive Frame
#define ISOTP_PCI_FC       0x30  // Flow Control
#define ISOTP_PADDING_BYTE   0xAA
#define ISOTP_FC_TIMEOUT_MS  1000
#define ISOTP_CF_SEP_MS      25
#define ISOTP_SF_MAX_PAYLOAD 7
#define ISOTP_FF_DATA_BYTES  6
#define ISOTP_CF_DATA_BYTES  7

// ---------- UDS (ISO 14229-1) ----------
#define UDS_SID_SESSION_CTRL     0x10  // DiagnosticSessionControl
#define UDS_SID_READ_DATA_BY_ID  0x22  // ReadDataByIdentifier
#define UDS_SESSION_DEFAULT      0x01
#define UDS_SESSION_EXTENDED     0x03

// ---------- HARDWARE ----------
#define ENGINE_START_PIN        5
#define IGNITION_HYSTERESIS_MS  300
#define DTC_FAULT_PIN           4
#define DTC_HYSTERESIS_MS       300

// ---------- NOISE & BROADCAST (opcionales, desactivados por defecto) ----------
#define BROADCAST_ENABLE       0   // 1 = emite tramas 0x280/0x320/0x3D0
#define BROADCAST_INTERVAL_MS  100
#define ADD_NOISE              0   // 1 = jitter en los sensores
\end{lstlisting}

\begin{lstlisting}[language=C++, caption={Funciones de codificación OBD-II — ecu\_protocol.h}]
// rpm*4 en dos bytes big-endian
inline void encodeRPM(uint8_t* data, uint16_t rpm) {
  uint16_t encoded = rpm * 4;
  data[0] = (encoded >> 8) & 0xFF;
  data[1] = encoded & 0xFF;
}

// Temperatura: T_raw = T_celsius + 40
inline uint8_t encodeTemp(int16_t temp) { return (uint8_t)(temp + 40); }

// Porcentaje: val_raw = percent * 255 / 100
inline uint8_t encodePercent(uint8_t percent) { return (percent * 255) / 100; }
\end{lstlisting}

\section*{A.2 — Interrupciones hardware e inyección de DTCs (\texttt{ecu\_sim.ino})}

Las dos interrupciones (ignición y fallo DTC) usan una ISR mínima con histéresis por
\texttt{millis()}; la lógica de efecto se ejecuta en el bucle principal.

\begin{lstlisting}[language=C++, caption={ISR con histéresis, inyección secuencial de DTCs y configuración}]
static const uint16_t DTC_FAULT_SET[]   = { DTC_P0300, DTC_P0171, DTC_P0420 };
static const uint8_t  DTC_FAULT_SET_SIZE = 3;

volatile uint16_t dtcFlag         = 0;
volatile uint32_t lastDtcISR      = 0;

// ISR minimalista: antirrebote por millis() + incremento de flag
void onDTCPress() {
  uint32_t now = millis();
  if ((uint32_t)(now - lastDtcISR) < DTC_HYSTERESIS_MS) return;
  lastDtcISR = now;
  dtcFlag++;
}

// Inyeccion secuencial; la 4a pulsacion (con 3 activas) borra todas
void handleDTCButton() {
  if (!dtcFlag) return;
  dtcFlag--;
  if (numStoredDTCs >= DTC_FAULT_SET_SIZE) {
    for (uint8_t i = 0; i < 8; i++) { dtcList[i].code = DTC_P0000; dtcList[i].active = false; }
    numStoredDTCs = 0; vehicle.numDTCs = 0; vehicle.checkEngine = false;
  } else {
    dtcList[numStoredDTCs].code   = DTC_FAULT_SET[numStoredDTCs];
    dtcList[numStoredDTCs].active = true;
    numStoredDTCs++;
    vehicle.numDTCs = numStoredDTCs; vehicle.checkEngine = true;
  }
}

void setup() {
  Serial.begin(SERIAL_BAUDRATE);
  CAN.begin(CAN_SPEED);
  pinMode(ENGINE_START_PIN, INPUT_PULLUP);
  attachInterrupt(digitalPinToInterrupt(ENGINE_START_PIN), onIgnitionPress, FALLING);
  pinMode(DTC_FAULT_PIN, INPUT_PULLUP);
  attachInterrupt(digitalPinToInterrupt(DTC_FAULT_PIN), onDTCPress, FALLING);
  initVehicleData();
  initDTCs();
}
\end{lstlisting}

\section*{A.3 — Recepción ISO-TP y despacho de modos}

\begin{lstlisting}[language=C++, caption={Parseo de Single Frame OBD-II y despacho a los modos/servicios}]
void processCANMessages() {
  if (CAN.parsePacket() <= 0) return;
  if (CAN.packetId() != ECU_CAN_ID) return;

  uint8_t data[8]; uint8_t dataLen = 0;
  while (CAN.available() && dataLen < 8) data[dataLen++] = (uint8_t)CAN.read();
  if (dataLen < 1) return;

  // Solo se procesan Single Frames de peticion
  if ((data[0] & 0xF0) != ISOTP_PCI_SF) return;

  uint8_t mode = data[1];
  uint8_t pid  = data[2];

  bool handled = false;
  switch (mode) {
    case MODE_01_CURRENT_DATA: handled = handleMode01(pid); break;
    case MODE_03_DTCS:         handled = handleMode03();    break;
    case MODE_04_CLEAR_DTCS:   handled = handleMode04();    break;
    case MODE_09_VEHICLE_INFO: handled = handleMode09(pid); break;
    case 0x10: handled = handleMode10(pid); break;            // UDS DSC
    case 0x22: {                                              // UDS RDBI
      uint8_t didLow = (dataLen >= 4) ? data[3] : 0x00;
      handled = handleMode22((uint16_t)(pid << 8) | didLow);
      break;
    }
    default:
      sendNegativeResponse(mode, NRC_SERVICE_NOT_SUPPORTED);
      handled = true;
      break;
  }
  if (!handled) sendNegativeResponse(mode, NRC_SUBFUNCTION_NOT_SUPP);
}
\end{lstlisting}

\section*{A.4 — Servicio UDS ReadDataByIdentifier (\texttt{0x22})}

\begin{lstlisting}[language=C++, caption={DIDs estándar y propietarios; control de acceso por sesión}]
bool handleMode22(uint16_t did) {
  uint8_t payload[24]; uint8_t payloadLen = 0;
  payload[0] = 0x62;                       // 0x22 + 0x40
  payload[1] = (did >> 8) & 0xFF;
  payload[2] = did & 0xFF;

  switch (did) {
    case 0xF190:                            // VIN (17 bytes, cualquier sesion)
      for (uint8_t i = 0; i < VIN_LENGTH; i++) payload[3 + i] = (uint8_t)vehicle.vin[i];
      payloadLen = 3 + VIN_LENGTH; break;
    case 0x2003:                            // RPM - solo sesion extendida
      if (udsSession != 3) { sendNegativeResponse(0x22, NRC_CONDITIONS_NOT_CORRECT); return true; }
      encodeRPM(&payload[3], vehicle.rpm); payloadLen = 5; break;
    // ... 0x2001, 0x2002, 0x2004..0x2008 analogos (extended only) ...
    default:
      sendNegativeResponse(0x22, NRC_REQUEST_OUT_OF_RANGE); return true;
  }

  if (payloadLen <= ISOTP_SF_MAX_PAYLOAD) {  // cabe en Single Frame
    memcpy(responseBuffer, payload, payloadLen);
    responseLength = payloadLen; sendResponse();
  } else {                                   // VIN: multi-frame
    sendMultiFrame(payload, payloadLen);
  }
  return true;
}
\end{lstlisting}

\section*{A.5 — Respuesta VIN con secuencia ISO-TP multi-frame (Modo \texttt{0x09})}

\begin{lstlisting}[language=C++, caption={First Frame + espera de Flow Control + Consecutive Frames}]
void sendVINMultiFrame() {
  const uint8_t PAYLOAD_LEN = 3 + VIN_LENGTH;          // 20 bytes
  uint8_t payload[20];
  payload[0] = MODE_09_VEHICLE_INFO + RESPONSE_SUCCESS; // 0x49
  payload[1] = 0x02; payload[2] = 0x01;
  for (uint8_t i = 0; i < VIN_LENGTH; i++) payload[3 + i] = (uint8_t)vehicle.vin[i];

  // First Frame: PCI 0x1L LL + primeros 6 bytes
  uint8_t ff[8];
  ff[0] = ISOTP_PCI_FF | ((PAYLOAD_LEN >> 8) & 0x0F);
  ff[1] = PAYLOAD_LEN & 0xFF;
  for (uint8_t i = 0; i < ISOTP_FF_DATA_BYTES; i++) ff[2 + i] = payload[i];
  CAN.beginPacket(ECU_RESPONSE_ID);
  for (uint8_t i = 0; i < 8; i++) CAN.write(ff[i]);
  CAN.endPacket();

  if (!waitForFlowControl()) Serial.println(F("[WARN] FC timeout"));

  // Consecutive Frames con SN incremental y separacion ISOTP_CF_SEP_MS
  uint8_t bytesSent = ISOTP_FF_DATA_BYTES, sn = 1;
  while (bytesSent < PAYLOAD_LEN) {
    uint8_t cf[8];
    cf[0] = ISOTP_PCI_CF | (sn & 0x0F);
    for (uint8_t i = 0; i < ISOTP_CF_DATA_BYTES; i++) {
      uint8_t idx = bytesSent + i;
      cf[1 + i] = (idx < PAYLOAD_LEN) ? payload[idx] : ISOTP_PADDING_BYTE;
    }
    CAN.beginPacket(ECU_RESPONSE_ID);
    for (uint8_t i = 0; i < 8; i++) CAN.write(cf[i]);
    CAN.endPacket();
    bytesSent += ISOTP_CF_DATA_BYTES;
    sn = (sn + 1) & 0x0F;
    if (bytesSent < PAYLOAD_LEN) delay(ISOTP_CF_SEP_MS);
  }
}
\end{lstlisting}
